\section{模型假设}
\label{sec:3}

\begin{enumerate}
\item \textbf{质量信号可作相对测量。} A 组指标经方向与域内尺度校正后，可比较同域文档的相对质量，不假定存在人工标注的绝对真值。该假设服务问题一的 $Q(x)$，并通过 A2/A3 冻结评分与删去极端指标的扰动检验其适用性。
\item \textbf{领域质量可按份额聚合。} 给定 A16 映射与桥接模型，混合语料质量取 17 域质量的份额加权值 $Q(\bp)=\sum_i p_iQ_i$；配比对验证损失的额外影响单独建模。该假设服务问题一至三，通过六个可映射域的桥接误差、A6/A7 留出与 A8–A11 检验其边界。
\item \textbf{评价目标预先固定。} 13 个验证域等权定义主目标，pile\_cc 单域仅作敏感性口径，评价权重不随候选配比改变。该假设服务问题一配比搜索与问题三目标选择。
\item \textbf{质量刻度需要锚定。} 主口径取 $Q_{\rm B}=Q(\bp)$，推荐配比处 $Q_0=Q(\bp^*)=0.6624116679870492$；$Q=1$ 且 $\bp=\bp_{\rm ref}$（A4 训练配比均值）时广义式退回经典式。副映射 $Q_{\rm B}=Q(\bp)/Q_0$ 会使部分领域超出 $[0,1]$，仅作敏感性处理。该假设服务问题二、三。
\item \textbf{成本按同一训练量计。} 训练、提质与注意力开销均以同一 $D$ 计量，成本函数参数依赛题给定，$\eta=2\times10^{-4}$ 为题设值。该假设服务问题三，并不预设三种成本函数下解的性质相同。
\item \textbf{非规模变化在窗口内平滑。} 问题四用连续时间项描述控制规模状态后的剩余得分变化，并以不同窗口与模型回测；遇到离散架构突破时该条件预测不适用。
\end{enumerate}
